% Author: Writer agent · 2026-11-14
% Task W1: Chapter 2 (signals). Sources: research-notes.md, Ch. 2. Graph: book/figures/noa_week.pdf (FICTIONAL data).

\chapter{כל דילוג נספר}\label{ch:2}

\begin{scene}
07:52. האוטובוס תקוע בפקק, ונועה מחפשת שיר שיעיר אותה. מתחילה בלדה שקטה, אחת האהובות עליה. אחרי 8 שניות נועה מדלגת. גם על הבא, ועל זה שאחריו. רק בשיר הרביעי, משהו מהיר עם תופים, היא נרגעת ומשעינה את הראש על החלון.

„יופי", היא חושבת, „עכשיו האפליקציה רשמה ביומן שלה שאני שונאת בלדות". היא שוקלת לרגע לשלוח לאפליקציה הודעת הבהרה. ואז נזכרת שאין לה למי.
\end{scene}

\section{משוב מפורש ומשוב משתמע}

כשמישהו נותן לסרט ארבעה כוכבים, הוא אומר לנו במפורש מה הוא חושב. זה \term{משוב מפורש}{explicit feedback}. אבל ברוב הזמן אנחנו לא נותנים ציון לשום דבר, אנחנו פשוט מקשיבים. לכן מערכות יכולות ללמוד גם מ\term{משוב משתמע}{implicit feedback}: הן עוקבות באופן פסיבי אחרי ההתנהגות שלנו, כמו היסטוריית רכישות, הרגלי צפייה או גלישה. בניגוד לציון מפורש, אין כאן אמירה ישירה של המשתמש על מה שהוא אוהב או לא אוהב \cite{hu2008collaborative}.

גם במוזיקה זה כך. לפי \en{Spotify}, פעולות כמו חיפוש, האזנה, דילוג ושמירה לספרייה משפיעות על האופן שבו החברה מפרשת את הטעם של המשתמש \cite{spotify2026understanding}.

\section{דילוג הוא רמז, לא פסק דין}

הבעיה במשוב משתמע היא שהוא רועש. לפי הו, קורן ווולינסקי, נתונים כאלה מבטאים העדפה חיובית וגם שלילית, וכל אחת מהן בדרגת ביטחון שונה מאוד \cite{hu2008collaborative}. דילוג אחד יכול לומר „לא אהבתי", אבל גם „לא עכשיו" או „אני מחפשת משהו אחר".

חוקרים מ-\en{Spotify} כינו את התנהגות הדילוג אחת משתי צורות חשובות במיוחד של משוב משתמע. הצורה השנייה היא מעבר בין הקשרי האזנה \cite{brost2019music}. במאגר הנתונים הציבורי שפרסמו, כ-160 מיליון סשנים של האזנה, הדילוג אינו כן או לא. המאגר מבחין בין שיר שהושמע „לזמן קצר מאוד", „לזמן קצר", שיר שרובו הושמע ושיר שהושמע עד הסוף \cite{brost2019music}. עוד מצאו החוקרים שדילוגים בסשן קשורים לדילוגים קודמים באותו סשן \cite{brost2019music}. לכן כדאי לקרוא דילוג בהקשר של הרגע כולו, ולא רק כשיפוט על השיר.

\begin{figure}[t]
  \centering
  \includegraphics[width=\textwidth]{noa_week.pdf}
  \caption{השבוע של נועה: שירים שהושמעו (יותר מ-30 שניות, לפי הגדרה שלנו לאיור זה) ושירים שדולגו, לפי שעה ביום \fictional}
  \label{fig:noa-week}
\end{figure}

איור~\ref{fig:noa-week} מציג שבוע אחד של נועה, ובו נתונים בדיוניים \fictional. לצורך האיור קבענו הגדרה משלנו: שיר „הושמע" אם נועה שמעה ממנו יותר מ-30 שניות. זו אינה הגדרה של פלטפורמה כלשהי. במהלך השבוע התחילו 344 שירים. 245 מהם הושמעו ועל 99 נועה דילגה, כלומר 28.8\%.

הפיזור לאורך היום מעניין יותר. באוטובוס, בין 7 ל-9 בבוקר, נועה דילגה על 34 מתוך 82 שירים (41\%). בחדר הכושר, בין 17 ל-19, על 11 מתוך 57 (19\%). בלימודים, בין 22 לחצות, רק על 5 מתוך 45 (11\%). האם נועה שונאת את שירי הבוקר? ממש לא. באוטובוס היא פשוט מחפשת את השיר הנכון. מערכת שהייתה מסיקה מכל דילוג „היא לא אוהבת את זה" הייתה טועה כאן בגדול.

\begin{realitycheck}
\textbf{כמה שירים אנשים מדלגים באמת?} לא מצאנו מקור מאומת שנותן אחוז כללי של דילוגים, ולכן איננו מביאים מספר כזה. כל המספרים באיור~\ref{fig:noa-week} בדיוניים, ונועדו רק להמחיש את הרעיון.
\end{realitycheck}

\begin{inshort}
\begin{itemize}
  \item משוב מפורש הוא מה שאנחנו אומרים. משוב משתמע הוא מה שאנחנו עושים, ואין בו אמירה ישירה על הטעם שלנו.
  \item דילוג הוא אות רועש: הוא יכול לומר „לא אהבתי", אבל גם „לא עכשיו".
  \item לפי \en{Spotify}, חיפוש, האזנה, דילוג ושמירה משפיעים על האופן שבו היא מפרשת את הטעם שלנו.
\end{itemize}
\end{inshort}
